% 同一离散场景的三种分割输出；字符在灰度下保留类别和身份区别。
\begin{tikzpicture}[font=\figurefont\small,x=4.5mm,y=-4.5mm]
 \foreach \offset/\title in {0/语义分割,8/实例分割,16/全景分割}{
  \begin{scope}[xshift=\offset*4.5mm]
   \node[anchor=south] at (3,0) {\title};
   \fill[BookPaper] (0,0.4) rectangle (6,4.4);
   \draw[BookMuted,line width=0.7pt] (0,0.4) rectangle (6,4.4);
   \draw[BookRule,line width=0.5pt] (0,1.4)--(6,1.4);
   \ifnum\offset=8
    \node[text=BookMuted] at (3,3.6) {不标背景};
   \else
    \node[text=BookMuted] at (3,3.6) {道路};
    \node[text=BookMuted] at (3,0.9) {天空};
   \fi
   \filldraw[fill=BookTeal!18,draw=BookTeal,line width=0.8pt]
      (0.5,1.8) rectangle (2.8,2.8);
   \ifnum\offset=0
    \filldraw[fill=BookTeal!18,draw=BookTeal,line width=0.8pt]
      (3.2,1.8) rectangle (5.5,2.8);
    \node at (1.65,2.3) {车};
    \node at (4.35,2.3) {车};
   \else
    \filldraw[fill=BookAmber!16,draw=BookAmber,dashed,line width=0.8pt]
      (3.2,1.8) rectangle (5.5,2.8);
    \node at (1.65,2.3) {车1};
    \node at (4.35,2.3) {车2};
   \fi
  \end{scope}
 }
\end{tikzpicture}
